\documentclass[11pt]{article}
\usepackage{amsmath,amssymb,geometry}
\geometry{margin=1in}
\title{ROOT V145: Topological Flux, Twisted Dirac Zero Modes, and Anomaly Audit}
\author{Reproducible mathematical research note}
\date{2026-11-24}
\begin{document}
\maketitle
\section{Status and units}
Known mathematical structure, not a proof of four-force unification.
Set only $c=1$; keep $\hbar$ and $G$ dimensionful.
\section{Two-patch U(1) connection on $S^2$}
\begin{align}
 A_N&=\frac n2(1-\cos\theta)\,d\phi,&
 A_S&=-\frac n2(1+\cos\theta)\,d\phi,\\
 A_N-A_S&=n\,d\phi,&
 F&=dA=\frac n2\sin\theta\,d\theta\wedge d\phi.
\end{align}
Single-valued transition $\exp(in\phi)$ implies $n\in\mathbb Z$.
\begin{equation} \frac1{2\pi}\int_{S^2}F=n. \end{equation}
\section{Twisted Dirac index and zero modes}
For integer-normalized charge $q$,
\begin{equation} \operatorname{ind}D_q=\frac q{2\pi}\int_{S^2}F=qn.\end{equation}
For $qn>0$, $\ker D_+\cong H^0(\mathbb{CP}^1,\mathcal O(qn-1))$ has dimension $qn$.
Its basis has locally $1,z,\ldots,z^{qn-1}$.
For $n>0$, normalized section integrals obey
\begin{equation}
\Vert z^k\Vert^2=4\pi R^2 B(k+1,n-k)
=4\pi R^2\frac{k!(n-k-1)!}{n!},\quad 0\le k<n.
\end{equation}
For the round sphere and monopole, the nonzero Dirac eigenvalues are
\begin{equation}
\lambda_{k,\pm}=\pm\frac{\sqrt{k(k+|qn|)}}R,
\qquad d_{k,\pm}=|qn|+2k,\quad k\ge1.
\end{equation}
The four-dimensional mass parameter, with $c=1$, is $m_k=\hbar|\lambda_{k,\pm}|$.
For converting to SI energy units only, use $E_k=\hbar_{\rm SI}c_{\rm SI}|\lambda_{k,\pm}|$.
\section{Single-hypercharge-flux obstruction}
With the integer normalization $\widetilde Y=6Y$, a single common flux $n=3$ gives
\begin{equation}
(\operatorname{ind}D_{Q_L},\operatorname{ind}D_{u^c},\operatorname{ind}D_{d^c},
\operatorname{ind}D_L,\operatorname{ind}D_{e^c})=(3,-12,6,-9,18),
\end{equation}
not equal multiplicities of three for all representations. This is a conditional
no-go for the minimal one-flux hypercharge model, not for every compactification.
\section{Anomaly audit}
For an isolated unit-charge left Weyl fermion with multiplicity three,
\begin{equation}\mathcal A_{U(1)^3}=3,\qquad\mathcal A_{\mathrm{grav}^2 U(1)}=3.\end{equation}
One Standard Model generation instead obeys
\begin{align}
\mathcal A_{Y^3}&=6(1/6)^3+3(-2/3)^3+3(1/3)^3+2(-1/2)^3+1=0,\\
\mathcal A_{\mathrm{grav}^2Y}&=6(1/6)+3(-2/3)+3(1/3)+2(-1/2)+1=0,\\
\mathcal A_{SU(2)^2Y}&\propto3(1/6)-1/2=0,\\
\mathcal A_{SU(3)^2Y}&\propto2(1/6)-2/3+1/3=0.
\end{align}
Additional six-dimensional anomaly cancellation and compactification stability are unverified.
\section{Unresolved physics}
Neither the round internal sphere nor its monopole flux follows from the proposed
space-light-speed helix. The radius, flux number, gauge representations, couplings,
fermion content, and gravitational origin remain inputs. Hence this does not
establish an original unified field theory.
\end{document}
